% Generated draft: requires mathematical and source review.
\documentclass[11pt]{article}
\usepackage[margin=25mm]{geometry}
\usepackage{fontspec}
\setmainfont{Latin Modern Roman}
\usepackage{amsmath,amssymb,mathtools,unicode-math}
\setmathfont{Latin Modern Math}
\usepackage{enumitem,needspace,xurl,hyperref}
\setlength{\parindent}{0pt}
\setlength{\parskip}{5pt}
\setlength{\emergencystretch}{3em}
\newcommand{\N}{\mathbb{N}}
\newcommand{\Z}{\mathbb{Z}}
\newcommand{\Q}{\mathbb{Q}}
\newcommand{\R}{\mathbb{R}}
\newcommand{\C}{\mathbb{C}}
\newcommand{\F}{\mathbb{F}}
\newcommand{\A}{\mathbb{A}}
\newcommand{\PP}{\mathbb{P}}
\newcommand{\E}{\mathbb{E}}
\newcommand{\eps}{\varepsilon}
\newcommand{\dd}{\,\mathrm{d}}
\begin{document}
\section*{UnsolvedMath ID: 30004702}
\textbf{Source problem code:} catalogue entry
\subsection{Definitions and mathematical statement}
The supplied record concerns the historical appropriation of mathematical classics by practitioners and craftsmen. It supplies no operational definitions of a \emph{mathematical classic}, \emph{appropriation}, \emph{transfer}, or \emph{good mathematics}, and specifies no corpus, historical period, geographical scope, or criteria for identifying the practitioners and craftsmen concerned.

\textbf{Source question.} When practitioners and craftsmen appropriated mathematical classics, what exactly was transferred: demonstrations, methods, practical results, ideas of good mathematics, or some combination of these? Can the works referred to in this discussion qualify as classics?

The listed forms of transfer are possibilities raised by the source; they are not specified as exhaustive or mutually exclusive. The referent of ``such works'' in the supplied question is unclear: it may concern the practitioners' and craftsmen's resulting works or the works they appropriated. Criteria for qualification as a classic are also unspecified. Accordingly, this is a historical research question whose precise scope and evaluative criteria remain unresolved, rather than a mathematical proposition with supplied definitions, domains, and quantifiers.
\subsection{Short English statement}
Investigate what mathematical content, practices, or standards practitioners and craftsmen acquired through appropriating mathematical classics, and whether the works concerned can themselves count as classics. The relevant works, historical scope, and meaning of classic require clarification.
\subsection{Sources}
\begin{itemize}
\item Local UnsolvedMath registry, supplied record with permanent ID \(30004702\), source problem code \texttt{catalogue entry}, entitled \emph{Mathematical Knowledge Transferred to Practitioners and Craftsmen}; collection: Oberwolfach Reports Open Problems. No individual proposer is supplied; the recorded proposed year is 2021.
\item Supplied report citation: \emph{Mathematics and its Ancient Classics Worldwide: Translations, Appropriations, Reconstructions, Roles} (2022), Oberwolfach Reports, pp.~1347--1406, DOI \texttt{10.4171/owr/2021/26}. Source URL and supplied literature source labeled ``Oberwolfach report'': \url{https://doi.org/10.4171/owr/2021/26}. No passage within this page range is identified.
\item Supplied literature assessment, checked \textbf{2026-08-22}: the transfer of mathematical content or practice and the classification of practitioners' works as classics remain open historical questions in the cited discussion. The record labels the problem \texttt{open}. This is supplied dated curation, not an independently verified assessment; time-sensitive status and cited claims require independent verification.
\item The record labels the statement \texttt{corrected\_verified} and reports manual semantic, typographic, and LaTeX review for self-containment and preservation of the mathematical target. These are supplied editorial claims, not independent verification of this draft or evidence that the historical question has a precise mathematical formulation.
\end{itemize}
\end{document}
